% Source: decision_tree.tex
\begin{notechapter}
\chapter{Decision Trees}\label{ch:decision-trees}

\section{Fundamentals}

\subsection{Overview}
A decision tree is a supervised machine learning algorithm used
for both classification and regression tasks. It works by
recursively splitting the data into subsets based on feature
values, aiming to maximize the homogeneity\footnote{Homogeneity
refers to the purity of a subset with respect to the target
variable. In classification, a subset is homogeneous if all
elements belong to the same class. In regression, a subset is
homogeneous if the target values are very similar.} of the target
variable within each subset.

\subsection{Key Concepts}

\subsubsection{Tree Structure}
\begin{itemize}
    \item \textbf{Nodes}: Represent decisions based on features.
    \item \textbf{Branches}: Represent outcomes of decisions.
    \item \textbf{Leaf Nodes}: Represent final predictions (class
    labels or continuous values).
\end{itemize}

\subsubsection{Splitting Criteria}
The choice of splitting criterion depends on the task:
\begin{itemize}
    \item \textbf{Classification}: Gini Impurity, Entropy,
    Information Gain, Gain Ratio.
    \item \textbf{Regression}: Variance Reduction.
\end{itemize}

\subsubsection{Recursive Partitioning}
The tree is built recursively by splitting the data at each node
until a stopping criterion is met (e.g., maximum depth, minimum
samples per leaf).

\subsubsection{Prediction}
To predict a new instance, traverse the tree from the root to a
leaf node based on feature values. The prediction is the value
associated with the leaf.

\section{Splitting Criteria}

\subsection{Gini Impurity}
Gini Impurity measures the likelihood of misclassifying a randomly
chosen element. It is calculated as:

\[
\text{Gini}(S) = 1 - \sum_{i=1}^{c} (p_i)^2
\]

where \( S \) is the dataset, \( c \) is the number of classes,
and \( p_i \) is the proportion of class \( i \) in \( S \).

\subsection{Entropy and Information Gain}
Entropy measures the disorder or uncertainty in the dataset:

\[
\text{Entropy}(S) = -\sum_{i=1}^{c} p_i \log_2(p_i)
\]

\subsubsection{Information Gain}
Information Gain measures the reduction in entropy after splitting
the dataset on a feature:

\[
\text{IG}(S, A) = \text{Entropy}(S) - \sum_{v \in \text{Values}(A)}
\frac{|S_v|}{|S|} \text{Entropy}(S_v)
\]

where \( A \) is the feature, \( S_v \) is the subset of \( S \)
where \( A = v \), and \( \text{Values}(A) \) are the possible
values of \( A \).

\subsubsection{Gain Ratio}
Gain Ratio is a modification of Information Gain that reduces bias
towards features with many values. It is calculated as:

\[
\text{GainRatio}(S, A) = \frac{\text{IG}(S, A)}{\text{SplitInfo}(A)}
\]

where \(\text{SplitInfo}(A)\) is the entropy of the feature \( A \):

\[
\text{SplitInfo}(A) = -\sum_{v \in \text{Values}(A)} \frac{|S_v|}{|S|}
\log_2 \left( \frac{|S_v|}{|S|} \right)
\]

\subsubsection{Why Split Info?}
Split Info penalizes features with many values by dividing the
Information Gain by the entropy of the feature itself. This
prevents the algorithm from favoring features that split the data
into many small subsets, which may not generalize well.

\subsection{Variance Reduction}
For regression tasks, Variance Reduction measures the reduction in
variance of the target variable after splitting:

\[
\text{Variance Reduction} = \text{Var}(S) - \sum_{v \in
\text{Values}(A)} \frac{|S_v|}{|S|} \text{Var}(S_v)
\]

\subsection{Misclassification Impurity}
Misclassification Impurity measures the minimum probability of
misclassifying a randomly chosen element:

\[
\text{Misclassification}(S) = 1 - \max(p_i)
\]

\section{Feature Importance in Decision Trees}
Feature importance is a measure of how much a feature contributes
to the predictive power of the decision tree. It is calculated
based on the reduction in the impurity measure (e.g., Gini Impurity
or Entropy for classification, Variance Reduction for regression)
that a feature achieves when it is used to split the data.

\subsubsection{How Feature Importance is Calculated}
\begin{itemize}
    \item \textbf{Impurity Reduction at Each Split}:
    \begin{itemize}
        \item At each node, the algorithm calculates the reduction
              in impurity achieved by splitting on a feature.
        \item The reduction is calculated as:
        \[
        \text{Impurity Reduction} = \text{Impurity}_{\text{before split}}
        - \left( \frac{|S_1|}{|S|} \cdot \text{Impurity}(S_1) + \frac{|S_2|}{|S|} \cdot \text{Impurity}(S_2) \right)
        \]
        where:
        \begin{itemize}
            \item \( S \) is the dataset before the split.
            \item \( S_1 \) and \( S_2 \) are the datasets in the left and right branches after the split.
            \item \( \text{Impurity}(S_1) \) and \( \text{Impurity}(S_2) \) are the impurities of the left and right branches.
            \item \( \frac{|S_1|}{|S|} \) and \( \frac{|S_2|}{|S|} \) are the proportions of data in the left and right branches.
        \end{itemize}
    \end{itemize}
    \item \textbf{Summing Impurity Reduction}:
    \begin{itemize}
        \item For each feature, the total impurity reduction is
              summed across all splits where the feature is used.
    \end{itemize}
    \item \textbf{Normalization}:
    \begin{itemize}
        \item The importance of each feature is normalized so that
              the sum of all feature importances is 1.
        \item The formula for feature importance is:
        \[
        \text{Importance}(A) = \frac{\sum_{\text{splits using } A}
        \text{Impurity Reduction}}{\sum_{\text{all splits}}
        \text{Impurity Reduction}}
        \]
    \end{itemize}
\end{itemize}

\subsubsection{Feature Usage in ID3 vs. CART}
\begin{itemize}
    \item \textbf{ID3}:
    \begin{itemize}
        \item In ID3, a feature can only be used \textbf{once in a given path}
              from the root to a leaf node.
        \item For example, if the feature "color" is used to split the data
              into branches for "blue," "red," and "green," the feature "color"
              cannot be reused in any of these branches.
        \item However, the same feature can be used in \textbf{different paths}
              (branches) of the tree, as long as it hasn't been used in the
              current path.
    \end{itemize}
    \item \textbf{CART}:
    \begin{itemize}
        \item In CART, a feature can be used \textbf{multiple times} in the tree,
              even in the same path, as long as it provides the best split.
        \item This allows CART to handle continuous features more effectively
              and create more flexible trees.
    \end{itemize}
\end{itemize}

\subsubsection{Example: Feature Usage in ID3}
Consider the following dataset:

\begin{tabular}{|l|l|l|}
    \hline
    Color & Size & Target \\
    \hline
    Blue & Small & Yes \\
    Blue & Large & No \\
    Red & Small & Yes \\
    Red & Large & No \\
    Green & Small & No \\
    Green & Large & Yes \\
    \hline
\end{tabular}

\subsubsection{Tree Construction}
\begin{itemize}
    \item \textbf{Root Node}:
    \begin{itemize}
        \item The algorithm selects "color" as the best feature to split on.
        \item The data is split into three branches: "blue," "red," and "green."
    \end{itemize}
    \item \textbf{Left Branch (Color = Blue)}:
    \begin{itemize}
        \item The subset of data for "blue" is:
        \[
        \text{Blue}: \{(Blue, Small, Yes), (Blue, Large, No)\}
        \]
        \item The feature "color" cannot be used again in this branch.
        \item The algorithm evaluates the remaining feature ("size") and splits the data based on it.
    \end{itemize}
    \item \textbf{Middle Branch (Color = Red)}:
    \begin{itemize}
        \item The subset of data for "red" is:
        \[
        \text{Red}: \{(Red, Small, Yes), (Red, Large, No)\}
        \]
        \item The feature "color" cannot be used again in this branch.
        \item The algorithm evaluates the remaining feature ("size") and splits the data based on it.
    \end{itemize}
    \item \textbf{Right Branch (Color = Green)}:
    \begin{itemize}
        \item The subset of data for "green" is:
        \[
        \text{Green}: \{(Green, Small, No), (Green, Large, Yes)\}
        \]
        \item The feature "color" cannot be used again in this branch.
        \item The algorithm evaluates the remaining feature ("size") and splits the data based on it.
    \end{itemize}
\end{itemize}

\subsubsection{Example: Calculating Feature Importance in CART}
Consider a decision tree with two features: Feature A and Feature B.
The tree has the following splits:

\begin{itemize}
    \item \textbf{Root Node}:
    \begin{itemize}
        \item Impurity before split: 0.5.
        \item Split on Feature A divides the data into two branches:
        \begin{itemize}
            \item Left branch: 60\% of the data, impurity = 0.4.
            \item Right branch: 40\% of the data, impurity = 0.2.
        \end{itemize}
        \item Impurity reduction:
        \[
        0.5 - \left( 0.6 \cdot 0.4 + 0.4 \cdot 0.2 \right) = 0.18
        \]
    \end{itemize}
    \item \textbf{Left Child Node}:
    \begin{itemize}
        \item Impurity before split: 0.4.
        \item Split on Feature B divides the data into two branches:
        \begin{itemize}
            \item Left branch: 50\% of the data, impurity = 0.3.
            \item Right branch: 50\% of the data, impurity = 0.1.
        \end{itemize}
        \item Impurity reduction:
        \[
        0.4 - \left( 0.5 \cdot 0.3 + 0.5 \cdot 0.1 \right) = 0.2
        \]
    \end{itemize}
    \item \textbf{Right Child Node}:
    \begin{itemize}
        \item Impurity before split: 0.2.
        \item Split on Feature A divides the data into two branches:
        \begin{itemize}
            \item Left branch: 70\% of the data, impurity = 0.1.
            \item Right branch: 30\% of the data, impurity = 0.3.
        \end{itemize}
        \item Impurity reduction:
        \[
        0.2 - \left( 0.7 \cdot 0.1 + 0.3 \cdot 0.3 \right) = 0.04
        \]
    \end{itemize}
\end{itemize}

\subsubsection{Calculating Total Impurity Reduction}
\begin{itemize}
    \item \textbf{Feature A}:
    \begin{itemize}
        \item Used in two splits with impurity reductions of 0.18 and 0.04.
        \item Total impurity reduction: \( 0.18 + 0.04 = 0.22 \).
    \end{itemize}
    \item \textbf{Feature B}:
    \begin{itemize}
        \item Used in one split with an impurity reduction of 0.2.
        \item Total impurity reduction: \( 0.2 \).
    \end{itemize}
\end{itemize}

\subsubsection{Calculating Feature Importance}
\begin{itemize}
    \item Total impurity reduction across all splits: \( 0.22 + 0.2 = 0.42 \).
    \item Normalized feature importance:
    \begin{itemize}
        \item Importance of Feature A: \( \frac{0.22}{0.42} \approx 0.52 \).
        \item Importance of Feature B: \( \frac{0.2}{0.42} \approx 0.48 \).
    \end{itemize}
\end{itemize}

\subsubsection{Interpretation}
\begin{itemize}
        \item Feature A contributes approximately 52\% to the model's predictive power.
    \item Feature B contributes approximately 48\% to the model's predictive power.
\end{itemize}

\begin{remark}[Weighting by node size]
In CART and in scikit-learn (``mean decrease in impurity''), the impurity reduction of a split is
weighted by the fraction of training samples that reach the node, \( N_t/N \): a split deep in the
tree, which only affects a few samples, counts less. Here the left child receives 60\% of the data and
the right child 40\%, so the weighted reductions are
\[
\text{A}: 1 \cdot 0.18 + 0.4 \cdot 0.04 = 0.196, \qquad \text{B}: 0.6 \cdot 0.2 = 0.12,
\]
giving normalized importances \( 0.196/0.316 \approx 0.62 \) for A and \( 0.12/0.316 \approx 0.38 \) for B.
The unweighted computation above overestimates the importance of features used in small nodes.
Note also that impurity-based importances are computed on the training data and are biased towards
features with many possible split points; permutation importance (below) avoids these issues.
\end{remark}

\section{Algorithms for Decision Trees}

\subsection{ID3 (Iterative Dichotomiser 3)}
\begin{itemize}
    \item \textbf{Purpose}: Classification.
    \item \textbf{Splitting Criterion}: Information Gain.
    \item \textbf{Pre-Pruning Support}: No.
    \item \textbf{Post-Pruning Support}: No.
    \item \textbf{Limitations}:
    \begin{itemize}
        \item Cannot handle continuous features.
        \item Prone to overfitting.
        \item No pruning support.
    \end{itemize}
\end{itemize}

\subsubsection{Algorithm Steps}
\begin{enumerate}
    \item Start with the entire dataset \( S \).
    \item Calculate the entropy of \( S \).
    \item For each feature \( A \) not yet used in the current path,
          calculate the Information Gain \( \text{IG}(S, A) \).
    \item Select the feature \( A \) with the highest Information Gain.
    \item Split \( S \) into subsets \( S_v \) based on the values of \( A \).
    \item Repeat steps 2--5 for each subset \( S_v \) until:
    \begin{itemize}
        \item All instances in a subset belong to the same class.
        \item No further splitting improves the purity of the subsets.
        \item All features have been used along the current path.
    \end{itemize}
    \item Assign a leaf node to each subset, labeled with the majority class.
\end{enumerate}

\subsection{C4.5}
\begin{itemize}
    \item \textbf{Purpose}: Classification.
    \item \textbf{Splitting Criterion}: Gain Ratio.
    \item \textbf{Pre-Pruning Support}: Yes.
    \item \textbf{Post-Pruning Support}: Yes.
    \item \textbf{Improvements over ID3}:
    \begin{itemize}
        \item Handles continuous features by discretization.
        \item Supports pruning (both pre-pruning and post-pruning).
        \item Handles missing values using probabilistic distributions.
        \item Supports boosting through instance weighting.
    \end{itemize}
\end{itemize}

\subsubsection{Algorithm Steps}
\begin{enumerate}
    \item Start with the entire dataset \( S \).
    \item Calculate the entropy of \( S \).
    \item For each feature \( A \) not yet used in the current path,
          calculate the Gain Ratio \( \text{GainRatio}(S, A) \).
    \item Select the feature \( A \) with the highest Gain Ratio.
    \item Split \( S \) into subsets \( S_v \) based on the values of \( A \).
    \item Repeat steps 2--5 for each subset \( S_v \) until:
    \begin{itemize}
        \item All instances in a subset belong to the same class.
        \item No further splitting improves the purity of the subsets.
        \item All features have been used along the current path.
    \end{itemize}
    \item Prune the tree to remove branches that do not improve generalization.
    \item Assign a leaf node to each subset, labeled with the majority class.
\end{enumerate}

\subsection{C5.0}
\begin{itemize}
    \item \textbf{Purpose}: Classification.
    \item \textbf{Splitting Criterion}: Gain Ratio.
    \item \textbf{Pre-Pruning Support}: Yes.
    \item \textbf{Post-Pruning Support}: Yes.
    \item \textbf{Improvements over C4.5}:
    \begin{itemize}
        \item More memory-efficient and faster.
        \item Better handling of noisy data.
    \end{itemize}
\end{itemize}

\subsubsection{Algorithm Steps}
\begin{enumerate}
    \item Start with the entire dataset \( S \).
    \item Calculate the entropy of \( S \).
    \item For each feature \( A \), calculate the Gain Ratio
          \( \text{GainRatio}(S, A) \).
    \item Select the feature \( A \) with the highest Gain Ratio.
    \item Split \( S \) into subsets \( S_v \) based on the
          values of \( A \).
    \item Repeat steps 2--5 for each subset \( S_v \) until:
    \begin{itemize}
        \item All instances in a subset belong to the same class.
        \item No further splitting improves the purity of the
              subsets.
    \end{itemize}
    \item Prune the tree to remove branches that do not improve
          generalization.
    \item Assign a leaf node to each subset, labeled with the
          majority class.
    \item (Optional) Apply boosting to improve accuracy.
\end{enumerate}

\subsection{CART (Classification and Regression Trees)}
\begin{itemize}
    \item \textbf{Purpose}: Classification and Regression.
    \item \textbf{Splitting Criterion}:
    \begin{itemize}
        \item Gini Impurity (Classification).
        \item Variance Reduction (Regression).
    \end{itemize}
    \item \textbf{Pre-Pruning Support}: Yes.
    \item \textbf{Post-Pruning Support}: Yes.
    \item \textbf{Tree Structure}: Binary splits.
    \item \textbf{Pruning}: Supports cost-complexity pruning.
    \item \textbf{Handling Missing Values}: Uses surrogate splits.
\end{itemize}

\subsubsection{Handling Categorical Features in CART}
For categorical features, CART handles them differently than
numerical features. Since categorical values (e.g., colors like
'red', 'blue', 'green') do not have a natural order, sorting is
not applicable. Instead, CART considers all possible binary
partitions of the categories and selects the one that minimizes
the impurity measure (e.g., Gini Impurity for classification or
Variance Reduction for regression).

\subsubsection{Binary Splits for Categorical Features}
CART always performs binary splits, even for categorical features.
For a feature with \( k \) possible values, the algorithm
considers all possible ways to split the \( k \) values into two
subsets. For example, if a feature has three categories: 'red',
'blue', and 'green', the possible binary splits would be:
\begin{itemize}
    \item \{red\} vs. \{blue, green\}
    \item \{blue\} vs. \{red, green\}
    \item \{green\} vs. \{red, blue\}
\end{itemize}

\subsubsection{Evaluating Splits}
For each possible binary split, CART evaluates the impurity
measure to determine how well the split separates the target
variable. The algorithm selects the binary split that maximizes
the homogeneity of the resulting subsets (i.e., minimizes
impurity).

\subsubsection{No Sorting Needed}
Since categorical values don't have an order, there's no need to
sort them. Instead, the algorithm considers all possible ways to
partition the categories into two groups and evaluates each
partition based on the impurity measure.

\subsubsection{Example}
Suppose you have a categorical feature "Color" with three possible
values: 'red', 'blue', and 'green'. The target variable is binary
(e.g., 'Yes' or 'No'). CART would evaluate the following splits:
\begin{itemize}
    \item \textbf{Split 1}: \{red\} vs. \{blue, green\}
    \item \textbf{Split 2}: \{blue\} vs. \{red, green\}
    \item \textbf{Split 3}: \{green\} vs. \{red, blue\}
\end{itemize}

The algorithm would then choose the split that results in the
lowest weighted average Gini Impurity across the two subsets.

\subsubsection{Surrogate Splits}
CART handles missing values using a technique called \textbf{surrogate splits}:
\begin{itemize}
    \item \textbf{Primary Split}: For each node, CART selects the best
          feature and split point to maximize homogeneity in the subsets.
    \item \textbf{Surrogate Splits}: If the primary splitting feature
          is missing for an instance, CART uses alternative features
          (called surrogate features) to determine the split.
    \item \textbf{How Surrogate Splits Are Learned}:
    \begin{itemize}
        \item During \textbf{training}, for each node, CART identifies
              the best surrogate features that most closely mimic the
              primary split.
        \item The surrogate features are ranked based on their ability
              to replicate the primary split.
        \item The best surrogate split is the one that most closely
              matches the primary split in terms of how it divides the
              data.
    \end{itemize}
    \item \textbf{When Are Surrogate Splits Used?}:
    \begin{itemize}
        \item \textbf{Training}: If the primary feature is missing for
              an instance during training, the instance is sent down
              the branch determined by the best surrogate split.
        \item \textbf{Testing}: During prediction, if the primary
              splitting feature is missing for an instance, the
              surrogate splits are used to determine the appropriate
              branch.
    \end{itemize}
    \item \textbf{Advantages of Surrogate Splits}:
    \begin{itemize}
        \item Allows CART to handle missing values without discarding
              instances or imputing values.
        \item Improves the robustness of the tree by providing
              alternative splitting rules.
    \end{itemize}
\end{itemize}

\subsubsection{Cost-Complexity Pruning}
CART uses \textbf{cost-complexity pruning} to prevent overfitting. The process involves:
\begin{itemize}
    \item \textbf{Cost-Complexity Parameter (\(\alpha\))}:
    \begin{itemize}
        \item \(\alpha\) controls the trade-off between the tree's
              complexity (number of leaves) and its accuracy.
        \item A higher \(\alpha\) results in a simpler tree with fewer
              leaves, while a lower \(\alpha\) allows for a more
              complex tree.
    \end{itemize}
    \item \textbf{Pruning Process}:
    \begin{itemize}
        \item If \(\alpha\) is fixed (e.g., \(\alpha = 0.25\)), the
              algorithm prunes the tree \textbf{once} to minimize the
              cost-complexity measure:
        \[
        R_{\alpha}(T) = R(T) + \alpha |T|
        \]
        where:
        \begin{itemize}
            \item \( R(T) \) is the error rate of the tree \( T \).
            \item \( |T| \) is the number of leaves in the tree.
        \end{itemize}
        \item If a \textbf{grid search} is performed, the algorithm evaluates
              a \textbf{sequence of subtrees}, each corresponding to a
              different value of \(\alpha\).
    \end{itemize}
    \item \textbf{Selecting the Optimal \(\alpha\)}:
    \begin{itemize}
        \item If \(\alpha\) is fixed, no further selection is needed.
        \item If a grid search is performed, the optimal \(\alpha\) can
              be chosen using \textbf{cross-validation} or based on a
              validation set.
    \end{itemize}
\end{itemize}

\subsubsection{Feature Selection in CART}
The CART algorithm selects the best feature to split on by evaluating
the impurity reduction achieved by each feature. The process involves
the following steps:

\begin{itemize}
    \item \textbf{Compute the Best Split for Each Feature}:
    \begin{itemize}
        \item For each feature, evaluate all possible splits and
              compute the impurity measure (e.g., Gini Impurity or
              Entropy) for each split.
        \item If using \textbf{Gini Impurity}, the best split is the one
              that minimizes:
        \[
        \text{Gini}_{\text{split}} = \frac{|S_1|}{|S|} \text{Gini}(S_1)
        + \frac{|S_2|}{|S|} \text{Gini}(S_2)
        \]
        \item If using \textbf{Entropy}, the best split is the one that
              maximizes \textbf{Information Gain}:
        \[
        \text{IG}(S, A) = \text{Entropy}(S)
        - \sum_{v \in \text{Values}(A)} \frac{|S_v|}{|S|} \text{Entropy}(S_v)
        \]
    \end{itemize}
    \item \textbf{Compute the Impurity Reduction}:
    \begin{itemize}
        \item If using \textbf{Gini Impurity}, the impurity reduction is:
        \[
        \text{Impurity Reduction} = \text{Gini}_{\text{before split}}
        - \text{Gini}_{\text{split}}
        \]
        \item If using \textbf{Entropy}, the impurity reduction is the
              \textbf{Information Gain}:
        \[
        \text{Impurity Reduction} = \text{IG}(S, A)
        \]
    \end{itemize}
    \item \textbf{Select the Best Feature}:
    \begin{itemize}
        \item The feature that provides the largest impurity reduction
              (or highest Information Gain) is selected for the split.
    \end{itemize}
\end{itemize}

\subsubsection{Example: Feature Selection Using Entropy and Information Gain}
Consider a dataset with two features: \textbf{Feature A} and \textbf{Feature B}.
To select the best feature for a split:
\begin{itemize}
    \item \textbf{Step 1: Compute the Best Split for Each Feature}:
    \begin{itemize}
        \item For \textbf{Feature A}, the best split results in:
        \[
        \text{IG}(S, A) = 0.2
        \]
        \item For \textbf{Feature B}, the best split results in:
        \[
        \text{IG}(S, B) = 0.3
        \]
    \end{itemize}
    \item \textbf{Step 2: Compute the Impurity Reduction}:
    \begin{itemize}
        \item The impurity reduction for \textbf{Feature A} is:
        \[
        \text{Impurity Reduction}(A) = 0.2
        \]
        \item The impurity reduction for \textbf{Feature B} is:
        \[
        \text{Impurity Reduction}(B) = 0.3
        \]
    \end{itemize}
    \item \textbf{Step 3: Select the Best Feature}:
    \begin{itemize}
        \item Since \textbf{Feature B} provides a larger impurity reduction
              (0.3) compared to \textbf{Feature A} (0.2), \textbf{Feature B} is
              selected for the split.
    \end{itemize}
\end{itemize}

\subsubsection{Advantages of CART}
CART has several advantages that make it a popular choice:
\begin{itemize}
    \item \textbf{Handles Both Tasks}: CART can be used for both
          classification and regression tasks, making it versatile.
    \item \textbf{Binary Splits}: CART always performs binary splits,
          simplifying the tree structure and making it easier to
          interpret.
    \item \textbf{Robust to Outliers}: CART is less sensitive to
          outliers because it uses binary splits and focuses on local
          splits.
    \item \textbf{Handles Mixed Data}: CART can handle both numerical
          and categorical features without extensive preprocessing.
    \item \textbf{Supports Pruning}: CART supports cost-complexity
          pruning, which helps prevent overfitting.
    \item \textbf{Handles Missing Values}: CART uses surrogate splits
          to handle missing values effectively.
    \item \textbf{Interpretability}: The tree structure is easy to
          visualize and interpret, making it useful for explaining
          model decisions.
\end{itemize}

\subsubsection{Finding the Best Split Point}
CART uses a greedy algorithm to find the best split point for each
feature. The process is as follows:

\begin{itemize}
    \item \textbf{For Numerical Features}:
    \begin{itemize}
        \item For each feature \( A \), sort the unique values of \( A \).
        \item For each possible split point \( v \) (midpoint between
              consecutive values):
        \begin{itemize}
            \item Split the dataset into two subsets:
            \begin{itemize}
                \item \( S_1 \): Instances where \( A \leq v \).
                \item \( S_2 \): Instances where \( A > v \).
            \end{itemize}
            \item Calculate the Gini Impurity for each subset.
            \item Compute the weighted average Gini Impurity:
            \[
            \text{Gini}_{\text{split}} = \frac{|S_1|}{|S|} \text{Gini}(S_1)
            + \frac{|S_2|}{|S|} \text{Gini}(S_2)
            \]
            \item Select the split point \( v \) that minimizes
                  \( \text{Gini}_{\text{split}} \).
        \end{itemize}
    \end{itemize}

    \item \textbf{For Categorical Features}:
    \begin{itemize}
        \item For each feature \( A \), consider all possible binary
              partitions of the categories.
        \item For each possible binary split:
        \begin{itemize}
            \item Split the dataset into two subsets based on the
                  partition.
            \item Calculate the Gini Impurity for each subset.
            \item Compute the weighted average Gini Impurity:
            \[
            \text{Gini}_{\text{split}} = \frac{|S_1|}{|S|} \text{Gini}(S_1)
            + \frac{|S_2|}{|S|} \text{Gini}(S_2)
            \]
            \item Select the binary split that minimizes
                  \( \text{Gini}_{\text{split}} \).
        \end{itemize}
    \end{itemize}

    \item \textbf{For Regression (Variance Reduction)}:
    \begin{itemize}
        \item For each feature \( A \), sort the unique values of \( A \).
        \item For each possible split point \( v \) (midpoint between
              consecutive values):
        \begin{itemize}
            \item Split the dataset into two subsets:
            \begin{itemize}
                \item \( S_1 \): Instances where \( A \leq v \).
                \item \( S_2 \): Instances where \( A > v \).
            \end{itemize}
            \item Calculate the variance of the target variable for
                  each subset.
            \item Compute the weighted average variance:
            \[
            \text{Var}_{\text{split}} = \frac{|S_1|}{|S|} \text{Var}(S_1)
            + \frac{|S_2|}{|S|} \text{Var}(S_2)
            \]
            \item Select the split point \( v \) that minimizes
                  \( \text{Var}_{\text{split}} \).
        \end{itemize}
    \end{itemize}
\end{itemize}

\subsubsection{Algorithm Steps}
\begin{enumerate}
    \item Start with the entire dataset \( S \).
    \item For classification, calculate the Gini Impurity of \( S \).
          For regression, calculate the variance of \( S \).
    \item For each feature \( A \), find the best split point that
          minimizes:
    \begin{itemize}
        \item Gini Impurity (for classification).
        \item Variance (for regression).
    \end{itemize}
    \item Split \( S \) into two subsets \( S_1 \) and \( S_2 \) based
          on the best split.
    \item Repeat steps 2--4 for each subset \( S_1 \) and \( S_2 \) until:
    \begin{itemize}
        \item All instances in a subset belong to the same class
              (classification).
        \item The variance of the target variable is minimized
              (regression).
    \end{itemize}
    \item Prune the tree using cost-complexity pruning to remove
          branches that do not improve generalization.
    \item Assign a leaf node to each subset, labeled with:
    \begin{itemize}
        \item The majority class (classification).
        \item The mean value of the target variable (regression).
    \end{itemize}
\end{enumerate}

\section{General Properties of Decision Trees}
\subsection{Advantages and Disadvantages of Decision Trees}
Decision trees are widely used in machine learning due to their simplicity and interpretability. However, they also have some limitations. Below is a list of their key advantages and disadvantages:

\subsubsection{Advantages}
\begin{itemize}
    \item \textbf{Interpretability}: Decision trees are easy to visualize and interpret, making them useful for explaining model decisions.
    \item \textbf{Handles Mixed Data}: Decision trees can handle both numerical and categorical features without requiring extensive preprocessing.
    \item \textbf{Non-Parametric}: Decision trees make no assumptions about the underlying data distribution, making them flexible and suitable for a wide range of datasets.
    \item \textbf{Handles Missing Values}: Algorithms like CART use surrogate splits to handle missing values effectively.
    \item \textbf{Supports Both Tasks}: Decision trees can be used for both classification and regression tasks.
    \item \textbf{Feature Importance}: Decision trees provide a measure of feature importance, which can be useful for feature selection.
    \item \textbf{Robust to Outliers}: Decision trees are less sensitive to outliers because they use local splits rather than global optimization.
    \item \textbf{No Need for Scaling}: Decision trees do not require feature scaling, as they are based on splitting rules rather than distance metrics.
\end{itemize}

\subsubsection{Disadvantages}
\begin{itemize}
    \item \textbf{Overfitting}: Decision trees are prone to overfitting, especially when they are deep and complex. Pruning and setting constraints (e.g., maximum depth) can help mitigate this.
    \item \textbf{High Variance}: Small changes in the data can lead to significantly different trees, making decision trees unstable. Ensemble methods like Random Forests can address this.
    \item \textbf{Bias Towards Dominant Features}: Decision trees can be biased towards features with many levels or dominant features, which can skew the splits.
    \item \textbf{Limited Extrapolation}: Decision trees are not well-suited for extrapolation tasks (e.g., predicting values outside the range of the training data).
    \item \textbf{Computational Complexity}: Building an optimal decision tree is computationally expensive, especially for large datasets with many features.
    \item \textbf{Imbalanced Data}: Decision trees can struggle with imbalanced datasets, as they may favor the majority class.
    \item \textbf{Not Suitable for Linear Relationships}: Decision trees are not ideal for datasets with strong linear relationships, as they partition the data into rectangular regions.
\end{itemize}

\section{Pruning}
Pruning reduces the size of a decision tree to prevent overfitting.

\subsection{Pre-pruning}
Stops tree construction early by setting constraints (e.g.,
maximum depth, minimum samples per leaf).

\subsection{Post-pruning}
Constructs the full tree and then removes branches that provide
little predictive power. Common methods include:

\subsubsection{Cost-Complexity Pruning}
Balances tree complexity and accuracy by minimizing:

\[
R_{\alpha}(T) = R(T) + \alpha |T|
\]

where:
\begin{itemize}
    \item \( R(T) \) is the error rate of the tree \( T \).
    \item \( |T| \) is the number of leaves in the tree.
    \item \( \alpha \) is a complexity parameter.
\end{itemize}

\subsubsection{Error Rate \( R(T) \)}
The error rate \( R(T) \) is computed as the proportion of
misclassified instances in the training set:

\[
R(T) = \frac{\text{Number of misclassified instances}}{\text{Total
number of instances}}
\]

\section{Ensemble Methods}

\subsection{Bagging}
Bagging (Bootstrap Aggregating) is a general-purpose ensemble 
technique applicable to any machine learning algorithm. It works as follows:
\begin{itemize}
    \item \textbf{Bootstrap Sampling}: Each model is trained on a 
    random subset of the data (with replacement).
    \item \textbf{Independence}: Models are trained independently, 
    making bagging suitable for parallelization.
    \item \textbf{Aggregation}: Predictions are combined by averaging 
    (for regression) or voting (for classification).
\end{itemize}

\subsubsection{Key Features of Bagging}
\begin{itemize}
    \item \textbf{Reduces Variance}: Bagging reduces variance in 
    high-variance algorithms (e.g., decision trees) by averaging models.
    \item \textbf{Improves Generalization}: Combining multiple models 
    helps prevent overfitting and improves generalization.
    \item \textbf{Applicable to Any Algorithm}: While often used with 
    decision trees (e.g., Random Forests), bagging works with others 
    like linear regression, SVMs, and neural networks.
\end{itemize}

\subsection{Random Forest}
Random Forest extends bagging for decision trees, adding randomness 
to reduce variance and improve generalization. Key features include:
\begin{itemize}
    \item \textbf{Bootstrap Sampling}: Each tree is trained on a 
    random subset of the data (with replacement).
    \item \textbf{Feature Randomness}: At each split, a random subset 
    of features is considered, reducing tree correlation.
    \item \textbf{Aggregation}: Predictions are averaged (regression) 
    or voted on (classification).
\end{itemize}

\subsubsection{Why Random Forests Were Introduced}
Random Forests address two main limitations of bagging with decision trees:
\begin{itemize}
    \item \textbf{High Correlation Between Trees}: In bagging, trees 
    can become highly correlated if strong features dominate splits. 
    Random Forests reduce this by considering random feature subsets.
    \item \textbf{Overfitting}: Additional randomness further reduces 
    overfitting compared to bagging.
\end{itemize}

\subsubsection{Out-of-Bag (OOB) Error}
Out-of-Bag (OOB) error estimates ensemble model performance without 
a separate validation set. It works as follows:
\begin{itemize}
    \item During bootstrap sampling, some instances are not included 
    in a tree's training set. These are "out-of-bag" (OOB) samples.
    \item Each tree predicts the OOB samples it did not see during training.
    \item The OOB error is the average error of these predictions 
    across all trees.
\end{itemize}

The OOB error is a reliable estimate of generalization error and can be used for:
\begin{itemize}
    \item Model evaluation (e.g., selecting hyperparameters).
    \item Feature importance estimation.
\end{itemize}

\section{Permutation Feature Importance}
Permutation Feature Importance is a model-agnostic method for computing
feature importance. It measures the increase in the model's prediction
error after permuting the feature's values. If permuting a feature's
values leads to a significant increase in error, the feature is
considered important.

\subsubsection{How Permutation Feature Importance Works}
\begin{itemize}
    \item \textbf{Train the Model}:
    \begin{itemize}
        \item Train a model (e.g., a Random Forest) on the training data.
    \end{itemize}
    \item \textbf{Out-of-Bag (OOB) Samples}:
    \begin{itemize}
        \item For each tree, identify the OOB samples (samples not used
              in training that tree).
    \end{itemize}
    \item \textbf{Compute Baseline Error}:
    \begin{itemize}
        \item For each tree, compute the prediction error on its OOB
              samples. This is the baseline error.
    \end{itemize}
    \item \textbf{Permute Feature Values}:
    \begin{itemize}
        \item For each feature, randomly permute its values in the OOB
              samples.
    \end{itemize}
    \item \textbf{Compute Permuted Error}:
    \begin{itemize}
        \item Use the permuted OOB samples to make predictions with the
              tree and compute the prediction error. This is the permuted
              error.
    \end{itemize}
    \item \textbf{Calculate Feature Importance}:
    \begin{itemize}
        \item The importance of the feature is the difference between
              the permuted error and the baseline error, averaged over
              all trees in the forest:
        \[
        \text{Importance}(A) = \frac{1}{N} \sum_{i=1}^{N} (\text{Permuted Error}_i - \text{Baseline Error}_i)
        \]
        where \( N \) is the number of trees in the forest.
    \end{itemize}
\end{itemize}

\subsubsection{Advantages of Permutation Feature Importance}
\begin{itemize}
    \item \textbf{Model-Agnostic}: Can be applied to any model.
    \item \textbf{Handles Correlated Features}: Less biased by correlated
          features compared to traditional importance measures.
    \item \textbf{Intuitive Interpretation}: Directly measures how much
          the model's performance degrades when a feature's information
          is destroyed.
\end{itemize}

\subsubsection{Limitations of Permutation Feature Importance}
\begin{itemize}
    \item \textbf{Computationally Expensive}: Can be time-consuming for
          large datasets.
    \item \textbf{Sensitive to Feature Correlations}: May underestimate
          importance for highly correlated features.
    \item \textbf{Dependence on Model Performance}: Requires a well-trained
          model for reliable importance scores.
\end{itemize}

\subsubsection{Example: Permutation Feature Importance in Random Forests}
Consider a Random Forest model trained on a dataset. To compute the
importance of a feature \( A \):
\begin{itemize}
    \item For each tree, identify the OOB samples and compute the baseline
          error (e.g., 0.1).
    \item Permute the values of feature \( A \) in the OOB samples.
    \item Compute the permuted error (e.g., 0.15).
    \item The importance of feature \( A \) for this tree is:
    \[
    \text{Importance}(A) = 0.15 - 0.1 = 0.05
    \]
    \item Repeat this process for all trees and average the results to
          get the final importance score for feature \( A \).
\end{itemize}

\subsection{Boosting}
Boosting combines multiple weak learners (e.g., shallow trees) to 
create a strong learner. It works by:
\begin{itemize}
    \item Training weak learners sequentially.
    \item Giving more weight to instances misclassified by previous learners.
\end{itemize}

\subsubsection{Instance Weights}
Instance weights assign higher importance to certain instances during 
training. In boosting, misclassified instances get higher weights in 
subsequent iterations.

\subsubsection{Why Some Algorithms Don't Support Instance Weights}
Some algorithms, like ID3, do not natively support instance weights because:
\begin{itemize}
    \item They were not designed to handle weighted data.
    \item Their splitting criteria (e.g., Information Gain) do not 
    incorporate instance weights.
\end{itemize}

In contrast, algorithms like C4.5 and C5.0 support instance weights, 
making them suitable for boosting.

\section{Summary of Algorithms}

\begin{tabular}{llllll}
    \toprule
    Algorithm & Purpose & Splitting Criterion & Pre-Pruning & Post-Pruning \\
    \midrule
    ID3 & Classification & Information Gain & No & No \\
    C4.5 & Classification & Gain Ratio & Yes & Yes \\
    C5.0 & Classification & Gain Ratio & Yes & Yes \\
    CART & Class. \& Reg. & Gini, Var. Red. & Yes & Yes \\
    \bottomrule
\end{tabular}

\section{Conclusion}
Decision trees are versatile algorithms with various splitting
criteria and techniques to handle overfitting. Pruning, boosting,
bagging, and Random Forests are essential tools to improve their
performance. Understanding these concepts helps in selecting the
right algorithm for specific tasks.

\section{Annex: Detailed Derivations and Examples}

\subsection{Derivation of Gini Impurity}
The goal is to calculate the probability of misclassifying a
randomly chosen element from the dataset. This depends on:
\begin{itemize}
    \item The proportion of each class in the dataset (\( p_i \)).
    \item The probability of misclassifying an element from each
    class.
\end{itemize}

\subsubsection{Step-by-Step Derivation}
\begin{enumerate}
    \item \textbf{Probability of Correct Classification for Class \( i \)}:
    For a class \( i \), the probability of correctly classifying
    a randomly chosen element is \( p_i \).

    \item \textbf{Probability of Misclassification for Class \( i \)}:
    The probability of misclassifying an element from class \( i \)
    is \( 1 - p_i \).

    \item \textbf{Weighted Misclassification Probability}:
    The overall probability of misclassification is the weighted
    sum of the misclassification probabilities for each class,
    where the weights are the class proportions \( p_i \):

    \[
    \sum_{i=1}^{c} p_i (1 - p_i)
    \]

    \item \textbf{Connection to Gini Impurity}:
    Expanding the formula:

    \[
    \sum_{i=1}^{c} p_i (1 - p_i) = \sum_{i=1}^{c} p_i - \sum_{i=1}^{c}
    p_i^2 = 1 - \sum_{i=1}^{c} p_i^2
    \]

    This is the Gini Impurity formula:

    \[
    \text{Gini}(S) = 1 - \sum_{i=1}^{c} p_i^2
    \]
\end{enumerate}

\subsection{Example: Gini Impurity Calculation}
Consider a dataset with 3 classes:
\begin{itemize}
    \item Class 1: \( p_1 = 0.5 \)
    \item Class 2: \( p_2 = 0.3 \)
    \item Class 3: \( p_3 = 0.2 \)
\end{itemize}

The Gini Impurity is:

\[
1 - (0.5^2 + 0.3^2 + 0.2^2) = 1 - (0.25 + 0.09 + 0.04) = 0.62
\]

\subsection{Example: Information Gain Calculation}
Consider a dataset \( S \) with 10 instances, split into two
subsets \( S_1 \) and \( S_2 \) based on a feature \( A \):
\begin{itemize}
    \item \( S_1 \): 6 instances (4 Class 1, 2 Class 2).
    \item \( S_2 \): 4 instances (1 Class 1, 3 Class 2).
\end{itemize}

The Entropy of \( S \) is:

\[
\text{Entropy}(S) = -\left( \frac{5}{10} \log_2 \frac{5}{10} +
\frac{5}{10} \log_2 \frac{5}{10} \right) = 1
\]

The Entropy of \( S_1 \) is:

\[
\text{Entropy}(S_1) = -\left( \frac{4}{6} \log_2 \frac{4}{6} +
\frac{2}{6} \log_2 \frac{2}{6} \right) \approx 0.918
\]

The Entropy of \( S_2 \) is:

\[
\text{Entropy}(S_2) = -\left( \frac{1}{4} \log_2 \frac{1}{4} +
\frac{3}{4} \log_2 \frac{3}{4} \right) \approx 0.811
\]

The Information Gain is:

\[
\text{IG}(S, A) = 1 - \left( \frac{6}{10} \times 0.918 +
\frac{4}{10} \times 0.811 \right) \approx 0.125
\]

\subsection{Example: Gain Ratio Calculation}
Consider the same dataset \( S \) with 10 instances, split into
two subsets \( S_1 \) and \( S_2 \) based on a feature \( A \):
\begin{itemize}
    \item \( S_1 \): 6 instances (4 Class 1, 2 Class 2).
    \item \( S_2 \): 4 instances (1 Class 1, 3 Class 2).
\end{itemize}

The Information Gain \( \text{IG}(S, A) \) is approximately 0.125
(as calculated above).

The Split Info for feature \( A \) is:

\[
\text{Split Info}(A) = -\left( \frac{6}{10} \log_2 \frac{6}{10} +
\frac{4}{10} \log_2 \frac{4}{10} \right) \approx 0.971
\]

The Gain Ratio is:

\[
\text{Gain Ratio}(S, A) = \frac{\text{IG}(S, A)}{\text{Split Info}(A)}
= \frac{0.125}{0.971} \approx 0.129
\]

\subsection{Handling Missing Values in C4.5}
C4.5 handles missing values using probabilistic distributions.
When a feature value is missing, the instance is fractionally
assigned to all possible values of the feature, weighted by
their observed probabilities. This ensures no data is discarded.

\subsubsection{Example: Entropy Calculation with Missing Values}
Consider a dataset \( S \) with 10 instances, where feature \( A \)
has two values: \( v_1 \) and \( v_2 \). Suppose:
\begin{itemize}
    \item 6 instances have \( A = v_1 \) (5 Class 1, 1 Class 2),
    \item 3 instances have \( A = v_2 \) (2 Class 1, 1 Class 2),
    \item 1 instance has a missing value for \( A \) (Class 1).
\end{itemize}

The observed probabilities for \( A \) are:
\[
p(v_1) = \frac{6}{9} = 0.67, \quad p(v_2) = \frac{3}{9} = 0.33.
\]

The instance with the missing value is fractionally assigned:
\begin{itemize}
    \item 67\% to \( v_1 \) (Class 1),
    \item 33\% to \( v_2 \) (Class 1).
\end{itemize}

Now, compute the entropy for each subset:

1. \textbf{Subset \( S_{v_1} \)}:
\begin{itemize}
    \item Total weight: \( 6 + 0.67 = 6.67 \).
    \item Class distribution:
    \begin{itemize}
        \item Class 1: \( 5 + 0.67 = 5.67 \),
        \item Class 2: \( 1 + 0 = 1 \).
    \end{itemize}
    \item Class probabilities:
    \[
    p_1 = \frac{5.67}{6.67} \approx 0.85, \quad
    p_2 = \frac{1}{6.67} \approx 0.15.
    \]
    \item Entropy:
    \[
    \text{Entropy}(S_{v_1}) = -0.85 \log_2(0.85) -
    0.15 \log_2(0.15) \approx 0.61.
    \]
\end{itemize}

2. \textbf{Subset \( S_{v_2} \)}:
\begin{itemize}
    \item Total weight: \( 3 + 0.33 = 3.33 \).
    \item Class distribution:
    \begin{itemize}
        \item Class 1: \( 2 + 0.33 = 2.33 \),
        \item Class 2: \( 1 + 0 = 1 \).
    \end{itemize}
    \item Class probabilities:
    \[
    p_1 = \frac{2.33}{3.33} \approx 0.70, \quad
    p_2 = \frac{1}{3.33} \approx 0.30.
    \]
    \item Entropy:
    \[
    \text{Entropy}(S_{v_2}) = -0.70 \log_2(0.70) -
    0.30 \log_2(0.30) \approx 0.88.
    \]
\end{itemize}

3. \textbf{Weighted Average Entropy}:
\[
\text{Entropy}_{\text{split}} = \frac{6.67}{10} \times 0.61 +
\frac{3.33}{10} \times 0.88 \approx 0.69.
\]

\subsubsection{Why This Approach?}
This method ensures missing values are handled without discarding
data or making arbitrary assumptions. By fractionally assigning
instances, C4.5 maintains statistical validity and uses all
available information.
\end{notechapter}
